\section{Retaining ambient surfaces through ordinary compression}
\label{sec:seam-repair}

The saturated-bundle operator concerns the topology of a supplied cut.  A
different obstruction to a fast hierarchy algorithm is repeated replacement
of later surfaces after an earlier surface is compressed.  The hierarchy
algorithm of~\cite[Sections~6 and~9]{lackenby2026} uses precisely such a
replacement step.  This section gives an additive repair bound for a
restricted, explicitly stated situation: the cut-open surfaces are the
pieces of a certified family of surfaces in one fixed ambient manifold.
The result does not construct the required families or provide a new
unknot-recognition algorithm.  Its purpose is to identify geometric data
under which the choice of every later ambient surface can be retained.

\subsection{The additional gluing obligation}

An ordinary sequential hierarchy does not automatically supply simultaneous
ambient surfaces.  After an earlier cut, a later surface can have unrelated
boundary traces on the two copies of the cutting surface.  Changing that
earlier surface can expose attachments on its opposite side which were
absent from the side on which a compression disk was found.  A list of
homology vectors does not record this information.

We use a \emph{surface atlas} to mean a concrete PL cut-and-reglue
presentation, with the surface pieces labeled and cooriented, and with the
following additional conditions.  Traces bearing the same label agree
under every specified face identification.  Every internal boundary arc is
paired exactly once; unpaired boundary lies on the true boundary of the
ambient manifold.  Coorientations agree across the identifications.
Distinct patches of a single label have disjoint interiors.  At the
identified edges and corners, the link of each label is a single embedded
circle at an interior surface point, or a single properly embedded interval
at a boundary surface point, in the corresponding ambient link.

These are conditions on actual charts and gluing maps.  In an explicit
finite refinement they are finite incidence and link tests.  No verifier
for an arbitrary binary-compressed surface atlas is implemented here.

In the following, a \emph{cut-open face} is a connected part of a labeled
surface cut along its intersections with the other surfaces, realized on
one boundary side of a cut piece.  It can have positive genus or several
boundary components; it is not required to be a polygon.  Its interior
maps injectively into the corresponding ambient surface.  A \emph{seam}
is an actual trace of one of the other cutting surfaces.  Auxiliary
subdivision edges in the atlas are not seams.  Thus essentiality below is
tested in the whole cut-open face, even when it is authenticated using
many polygon charts.

\begin{lemma}[What an authenticated atlas supplies]
\label{lem:seam-atlas}
Under the preceding conditions, the quotient of the patches with any one
label is a properly embedded cooriented PL surface.  In an orientable
ambient manifold it is oriented.  Suppose an embedded disk in one cut
piece has interior disjoint from its boundary faces, and has boundary in
the interior of a single cut-open face of one label, disjoint from
all seams and other labels.  Its image under regluing is an embedded disk
whose interior avoids the whole surface family and whose boundary meets
only that labeled surface.
\end{lemma}

\begin{proof}
At a patch interior the required local surface chart is already supplied.
Across an identified face, the paired traces give two compatible local
halves of a surface.  The edge and corner link conditions supply the local
product charts at the remaining points; in particular they exclude
branching, self-intersection, and an unpaired internal boundary.  The
boundary condition identifies the resulting surface boundary with a subset
of the true manifold boundary.  The coorientations glue, and an ambient
orientation converts a coorientation into a surface orientation.

Regluing is injective on the interior of each cut piece.  On the boundary
of the stated disk, it is injective because this boundary lies inside one
cut-open face whose interior embeds in its labeled ambient surface.
No interior point of the
disk can be identified with a boundary point: the former is off the
surface family and the latter lies on it.  The disk image is therefore
embedded.  The same conditions show that it meets no other labeled
surface and meets its own label only along its boundary.
\end{proof}

The local link tests in Lemma~\ref{lem:seam-atlas} must not be omitted.
Agreement of formal boundary chains alone would allow a branched or
immersed quotient.  Conversely, once the ambient surfaces have been
authenticated, the repair theorem below is purely geometric and does not
depend on a particular triangulation of them.

\subsection{A potential for one fixed ambient family}

Let $M$ be a compact orientable irreducible $3$-manifold, and let $P$ be a
boundary pattern.  Let
\[
                         \mathcal F=(F_1,\ldots,F_k)
\]
be connected oriented surfaces properly embedded in $M$, in general
position and transverse to $P$.  Assume that the classes
$[F_1],\ldots,[F_k]$ are linearly independent in
$H_2(M,\partial M;\mathbb Q)$.  Pairwise intersections are finite unions
of arcs and circles.  Triple intersections, if present, are included in
the atlas; they do not split a component of a pairwise intersection.

Use the pattern complexity from~\cite[Section~3]{lackenby2026}, measured
against the fixed original pattern $P$:
\begin{equation}
 c_P(F_i)=-4\chi(F_i)+|\partial F_i\cap P|,
 \qquad
 \iota(\mathcal F)=
       \sum_{1\le i<j\le k}\#\pi_0(F_i\cap F_j).
 \label{eq:seam-complexities}
\end{equation}
Here $\#\pi_0$ counts geometric connected components, with no algebraic
cancellation.  The pattern in~\eqref{eq:seam-complexities} is not the
changing pattern of an intermediate cut.

\begin{lemma}[Negative complexity and a nonzero relative class]
\label{lem:seam-negative}
Let $F$ be a connected oriented surface properly embedded in $M$ with
$[F]\ne0$ in $H_2(M,\partial M;\mathbb Q)$.  If $c_P(F)<0$, then $F$
is a violating disk for $(M,P)$.
\end{lemma}

\begin{proof}
A connected orientable surface of negative pattern complexity is either
a sphere or a disk meeting $P$ at most three times.  An embedded sphere
bounds a ball by irreducibility, so its relative homology class is zero.
This contradicts the nonzero-class assumption.  If a disk is not violating,
the definition supplies a disk in
$\partial M$ with the same boundary and with the permitted empty, arc, or
tripod intersection with $P$.  Together the two disks bound a ball, after
the usual small boundary perturbation, by irreducibility.  The original
disk is boundary-parallel and represents zero in relative homology,
again contradicting the hypothesis.  Thus the disk is violating.
\end{proof}

Accordingly, suppose throughout a nonterminal stage that every
$c_P(F_i)$ is nonnegative.  Define
\begin{equation}
                  Q(\mathcal F)
                  =\sum_{i=1}^k c_P(F_i)+\iota(\mathcal F).
                  \label{eq:seam-potential}
\end{equation}
This is a nonnegative integer.  The intersection term is essential: a
loop can be essential in a cut-open face while bounding a disk in its
uncut ambient surface.  Such a loop leads to a disk exchange, rather
than to a reduction in the surface's Euler characteristic.

\begin{theorem}[Ordinary compression repair in an ambient block]
\label{thm:seam-repair}
Assume the preceding hypotheses and an authenticated atlas for
$\mathcal F$.  Suppose a disk $D$ in one of its cut pieces satisfies the
disk hypotheses of Lemma~\ref{lem:seam-atlas}, and that its boundary
$\alpha$ is essential in that cut-open face, which has label $j$.
Then one can either obtain a violating disk for $(M,P)$ or replace
$F_j$ by a connected oriented proper surface $G$, retaining all the
other ambient surfaces, such that their $k$ relative classes remain
independent and
\[
                   Q(F_1,\ldots,G,\ldots,F_k)
                      <Q(\mathcal F).
\]
More precisely:
\begin{enumerate}
\item If $\alpha$ is essential in the ambient surface $F_j$, then
      $G$ can be chosen among the components of the compression of
      $F_j$ along $D$, and in the nonterminal case
      \[
          c_P(G)\le c_P(F_j)-4,
          \qquad \iota(\mathcal F_{\mathrm{new}})
                         \le\iota(\mathcal F).
      \]
\item If $\alpha$ bounds a disk in the ambient $F_j$, then $G$ can
      be obtained by a disk exchange.  It has the same relative class,
      Euler characteristic, and boundary as $F_j$, while
      \[
                     \iota(\mathcal F_{\mathrm{new}})
                         <\iota(\mathcal F).
      \]
\end{enumerate}
\end{theorem}

\begin{proof}
By Lemma~\ref{lem:seam-atlas}, regluing gives an embedded disk in $M$
meeting $F_j$ exactly in its boundary and avoiding every $F_i$ with
$i\ne j$.  We continue to denote this image by $D$.  A sufficiently
small surgery neighborhood of $D$ is disjoint from the other surfaces.

First suppose that $\alpha$ is essential in $F_j$.  Let $F'_j$ be the
oriented surface obtained by compression.  The surgery trace gives
\begin{equation}
                   [F'_j]=[F_j]
                   \quad\hbox{in }H_2(M,\partial M;\mathbb Q).
                   \label{eq:seam-class}
\end{equation}
Write $V$ for the span of the other $k-1$ classes.  At least one
connected component $G$ of $F'_j$ has $[G]\notin V$, since otherwise
their sum~\eqref{eq:seam-class} would belong to $V$.  Choosing such a
component preserves rank $k$.  In particular $[G]\ne0$.  If
$c_P(G)<0$, Lemma~\ref{lem:seam-negative}, with the newly independent
family, supplies the asserted violating disk.

For the remaining case we prove the quantitative drop, including the
effect of a possibly discarded component.  Compression increases total
Euler characteristic by two and leaves the boundary and its intersections
with $P$ unchanged.  Consequently
\begin{equation}
                       c_P(F'_j)=c_P(F_j)-8.
                       \label{eq:seam-eight}
\end{equation}
There are at most two components.  No component is a sphere.  Indeed, a
separating compression would produce a sphere only if $\alpha$ already
bounded a disk in $F_j$, contrary to this case.  A nonseparating
compression produces a sphere only when $F_j$ is a closed torus.  Then
\eqref{eq:seam-class} and irreducibility would force $[F_j]=0$, again
contrary to independence.

Every connected sphere-free orientable surface has Euler characteristic
at most one, so a discarded component, if present, has pattern
complexity at least $-4$.  Subtracting that complexity
from~\eqref{eq:seam-eight} yields
$c_P(G)\le c_P(F_j)-4$.  All intersections with the other surfaces
are disjoint from the surgery neighborhood.  Each lies entirely in a
component of the compressed surface, because it misses the compressed
circle.  Selecting $G$ therefore deletes entire pairwise intersection
arcs or circles and introduces none.  This proves the first case.

Now suppose that $\alpha$ bounds a disk $E\subset F_j$.  It must
contain a pairwise intersection circle: if its interior missed every
other surface, it would lift to a disk in the same cut-open face,
contradicting the essentiality hypothesis there.  Since $\alpha$ is
disjoint from the other surfaces and $E$ lies in the interior of
$F_j$, every pairwise intersection component meeting $E$ is a whole
circle contained in its interior.

The embedded sphere $E\cup D$ bounds a ball $B$ by irreducibility.
The connected surface $F_j\setminus\operatorname{int}E$ lies outside
this ball.  If it lay inside, then in the closed-surface case the whole
$F_j$ could be pushed into $B$, forcing its relative class to be zero.
In the case with boundary, its boundary on $\partial M$ rules out that
possibility, since $B$ is bounded by a sphere in the interior of $M$.
Thus the usual disk exchange across $B$ replaces $E$ by a small
parallel copy of $D$ while retaining an embedded oriented proper
surface.  The isotopy may pass through other members of the family;
their mutual disjointness during the isotopy is not a hypothesis.

The replacement disk is disjoint from all the other surfaces.  It
therefore deletes every intersection circle in $E$ and creates no new
intersection.  At least one circle is deleted.  The surface's relative
class, Euler characteristic, and boundary are unchanged.  This proves
the second case, and hence the strict decrease of
\eqref{eq:seam-potential} in either nonterminal case.
\end{proof}

\begin{corollary}[Additive bound for the specified repairs]
\label{cor:seam-budget}
Starting from one family satisfying the theorem, any sequence of its
specified nonterminal repairs has length at most
\[
             \left\lfloor\frac{C_0}{4}\right\rfloor+I_0,
 \qquad C_0=\sum_i c_P(F_i^{\mathrm{initial}}),
 \qquad I_0=\iota(\mathcal F_{\mathrm{initial}}).
\]
This is at most $Q(\mathcal F_{\mathrm{initial}})$.  If the last repair
instead terminates by producing a violating disk, the total number of
repairs is at most
$\lfloor C_0/4\rfloor+I_0+1$.
\end{corollary}

\begin{proof}
The $k$ classes remain independent.  At each nonterminal stage their
pattern complexities are nonnegative.  Each nonterminal compression
along a curve essential in its ambient surface reduces their sum by at
least four; disk exchanges preserve that sum.  There can therefore be at
most $\lfloor C_0/4\rfloor$ such compressions.  Each disk exchange
decreases the intersection count by at least one, and a compression never
increases that count.  There can consequently be at most $I_0$ disk
exchanges.  Adding these distinct operation counts proves the claimed
bound.  A terminal surgery is allowed to leave the nonnegative-complexity
regime, so it is counted separately.
\end{proof}

\begin{corollary}[Retaining individually nonzero labels]
\label{cor:seam-nonzero}
The geometric repair and additive bound remain valid if mutual
independence is replaced by the weaker assumption
\[
               [F_i]\ne0\quad\text{for every }i.
\]
In this version every retained label has a connected representative with
nonzero relative class.  No mutual independence conclusion is asserted.
The number of labels is unrestricted by the dimension of the homology
group.  As before, the initial stage is nonterminal only if all its
pattern complexities are nonnegative.
\end{corollary}

\begin{proof}
For an essential compression, the oriented sum of the classes of its
components is the nonzero class $[F_j]$.  Hence at least one component
$G$ has nonzero class.  Choose that component.  The exclusion of sphere
components in the proof of Theorem~\ref{thm:seam-repair} uses only
$[F_j]\ne0$ and irreducibility.  The resulting estimate
$c_P(G)\leq c_P(F_j)-4$ is therefore unchanged.  If $c_P(G)<0$,
Lemma~\ref{lem:seam-negative} still produces a violating disk.
Otherwise the same intersection argument proves a strict decrease of
$Q$, and all labels still have nonzero classes.

In the disk-exchange case, the only homological input was that $F_j$
cannot be pushed into a ball.  Its nonzero class supplies exactly that
input.  The exchange preserves its individual class and strictly
decreases the intersection term.  The separate compression and
disk-exchange counts in Corollary~\ref{cor:seam-budget} are therefore
unchanged: at most $\lfloor C_0/4\rfloor+I_0$ nonterminal repairs
and at most one additional terminal repair.
\end{proof}

Keeping every component of $F'_j$ as one labeled oriented surface would
preserve the individual class $[F_j]$.  The connected-component selection
used in the theorem preserves independence and the number of labels, but
need not preserve that individual class.  This distinction matters if a
later construction requires one fixed cohomology basis rather than merely
an independent family.

\subsection{What the additive estimate does not establish}

\paragraph{Ambient retention still requires updated cut data.}
The unchanged objects are the other surfaces embedded in $M$.  Their
pieces in the successive cuts, the induced boundary patterns, and their
source certificates can change.  Discarding one compressed component can
glue together several old pieces of a later surface.  The theorem therefore
permits retaining ambient surface choices; it does not justify reusing
old normal coordinates, handle structures, source digests, or suffix
certificates without verification.

\paragraph{The disks have restricted boundary behavior.}
The hypotheses cover an ordinary compression whose boundary stays in one
surface face and misses every seam.  They do not cover general
pattern-compressions involving arcs on the original boundary or several
labeled faces.  Such an extension must account for the changed pattern
intersections and for the reconnection of seam traces.  Nor does the theorem
find all the disks it permits.

\paragraph{No simultaneous normalization bound is supplied.}
The repaired surfaces need not be normal in the previous handle structure.
Normalizing one of them can change its intersections with the others.  To
use Corollary~\ref{cor:seam-budget} inside a normal-surface algorithm, one
would need to control that change or charge it to another bounded
quantity.  Likewise, irreducibility proves that the disk-exchange ball
exists; this proof supplies no bound for a compressed description of the
resulting isotopy.

\paragraph{The potential is geometric, not automatically polynomial in
the input bit length.}
For example, take the two vertical annuli in $T^2\times I$ above curves
of slopes $(1,0)$ and $(1,N)$ in minimal position.  Their pattern
complexities for $P=\varnothing$ are zero, but their pairwise intersection
has $N$ arcs.  A slope or twist description can store $N$ in binary.
Thus even an additive bound in $Q$ can be exponential in a compressed
description length.  The example illustrates the size of the bound; it
does not assert that these annuli actually require $N$ repairs.

\paragraph{Homological independence has limited force.}
The preserved assertions are embeddedness, orientation, properness,
connectedness of each selected representative, and independence in the
fixed ambient relative homology group.  The theorem asserts neither
incompressibility nor pattern-incompressibility.  It also does not prove
nonseparation of every later cut-open surface after several previous cuts,
or the $q$-complexity descent required by the hierarchy algorithm.

For the independent version there is also a strict rank restriction:
$k\leq\dim H_2(M,\partial M;\mathbb Q)=b_1(M)$ by
Poincar\'e--Lefschetz duality.  For a knot exterior, abelianizing a
Wirtinger presentation identifies all arc generators with one meridian
generator, so $b_1(M)=1$.  The independent version consequently cannot
retain a whole multisurface hierarchy in the original knot exterior.
Its multiple-label use requires an appropriate later ambient block of
larger first Betti number and a suitable atlas there.
Corollary~\ref{cor:seam-nonzero} avoids this rank restriction by retaining
only individual nonzero classes, and so permits several ambient labels
even in a knot exterior.  Neither version proves that the resulting
pieces form a valid sequential hierarchy.  In particular, the corollary
does not replace the missing later-cut nonseparation argument.

The resulting research problem is concrete.  Construct ambient atlases
whose geometric intersection complexity is controlled, authenticate their
seams without expanding binary multiplicities, and retain that control
through simultaneous normalization and pattern-compression.  The local
cutting and bundle results, including the incoming relative-handle
work~\cite{incoming-hierarchy-kernels}, address different parts of this
construction.  The additive repair theorem does not close those remaining
obligations or imply a global quasi-polynomial recognition bound.
